% called by main.tex
%
\chapter{Introducción}
\label{ch::introduccion}

Este Trabajo Fin de Máster estudia un problema de predicción financiera de muy corto
plazo en un entorno poco habitual: los mercados de predicción sobre el precio de Bitcoin
de la plataforma Polymarket. El objetivo no es un sistema de trading en producción, sino
una \emph{metodología} rigurosa para responder, con honestidad estadística, a una
pregunta concreta: ¿contiene el libro de órdenes de estos mercados señales que permitan
anticipar el movimiento del precio en los próximos segundos y que, además, sobrevivan a
los costes y a la latencia reales de operar?

Como el lector puede no estar familiarizado con el mundo de las criptomonedas ni con la
microestructura de mercado, este capítulo introduce todos los conceptos desde cero antes
de plantear el problema y los objetivos.

\section{Introducción y definición del problema}

Esta sección define el problema desde cero. Antes de plantear la pregunta de
investigación se introducen los conceptos imprescindibles —mercado de predicción,
contrato binario, libro de órdenes y señales de microestructura—, de modo que el documento
resulte autocontenido también para un lector ajeno al mundo de las criptomonedas. Con ese
vocabulario fijado, se enuncia el problema concreto que se aborda.

\subsection{Qué es un mercado de predicción}

Un \emph{mercado de predicción} es un mercado en el que se compran y venden contratos
cuyo valor depende del resultado de un acontecimiento futuro. El caso más sencillo es el
\emph{contrato binario}: un contrato que pagará una unidad monetaria si el acontecimiento
ocurre y cero si no ocurre. Por construcción, su precio se mueve entre 0 y 1, y puede
leerse directamente como la \emph{probabilidad implícita} que el mercado asigna al
suceso. Por ejemplo, un contrato ``el precio de Bitcoin estará por encima de un umbral a
una hora determinada'' que cotiza a 0,60 indica que el mercado estima una probabilidad
cercana al 60\,\% de que eso ocurra. Cuando el acontecimiento se resuelve, el contrato
vale 1 o 0 y las posiciones se liquidan. La literatura clásica ha mostrado que estos
mercados agregan información dispersa y producen previsiones razonablemente
calibradas~\cite{wolfers2004prediction}.

\subsection{Qué es Polymarket}

Polymarket es una plataforma de mercados de predicción descentralizada. A diferencia de
una casa de apuestas, donde el operador fija las cuotas, en Polymarket son los propios
participantes quienes ponen precios mediante un \emph{libro de órdenes}: el emparejamiento
de compras y ventas ocurre fuera de la cadena de bloques (de forma rápida) y la
liquidación final del dinero se registra en cadena (sobre una \emph{blockchain})
\cite{polymarket_trading}. En los mercados sobre el precio de Bitcoin, cada contrato
binario toma valores entre 0 y 1 y se interpreta como una
probabilidad~\cite{polymarket_prices}. Esto convierte cada mercado en un pequeño
\emph{mercado electrónico} (\emph{exchange}) con su propia dinámica de oferta y demanda.

\subsection{El libro de órdenes y el vocabulario de la microestructura}

El \emph{libro de órdenes} (en inglés \emph{limit order book}, LOB) es la lista, ordenada
por precio, de todas las órdenes de compra y de venta que están a la espera de ejecutarse.
Sus elementos esenciales son:

\begin{itemize}
  \item \textbf{Bid}: el precio más alto al que alguien está dispuesto a comprar.
  \item \textbf{Ask}: el precio más bajo al que alguien está dispuesto a vender. Siempre
  es mayor o igual que el \emph{bid}.
  \item \textbf{Spread}: la diferencia entre el \emph{ask} y el \emph{bid}. Es el primer
  coste de operar: quien compra al \emph{ask} y vende al \emph{bid} pierde el \emph{spread}.
  \item \textbf{Profundidad}: el volumen disponible en cada nivel de precio. A más
  profundidad, más se puede operar sin mover el precio.
  \item \textbf{Precio medio} (\emph{mid}): el punto medio entre \emph{bid} y \emph{ask};
  una referencia ``limpia'' del precio sin el ruido del \emph{spread}.
  \item \textbf{Tick}: la variación mínima de precio que admite el mercado. Es la unidad
  natural en la que mediremos todos los resultados de este trabajo.
\end{itemize}

La \emph{microestructura de mercado} es la rama que estudia cómo se forman los precios a
partir de estos mecanismos —órdenes, \emph{spread}, profundidad y flujo de
operaciones— a escalas de tiempo muy finas. La Figura~\ref{fig:mercado} resume el entorno:
el precio del contrato se forma en su propio libro de órdenes, pero está acoplado a
referencias externas del mercado de Bitcoin (mercado al contado, futuros perpetuos y un
oráculo de precios), porque al fin y al cabo el contrato habla del precio de Bitcoin.

\begin{figure}[h]
\centering
\begin{tikzpicture}[
  node distance=8mm,
  box/.style={draw, rounded corners, align=center, minimum height=9mm, inner sep=4pt, font=\small},
  ext/.style={box, fill=blue!6},
  core/.style={box, fill=orange!12, thick},
  tgt/.style={box, fill=green!10},
  >=latex
]
\node[ext] (spot) {Mercado al contado\\(spot BTC)};
\node[ext, below=of spot] (perp) {Futuros perpetuos\\(perp BTC)};
\node[ext, below=of perp] (oracle) {Oráculo de precios};
\node[core, right=22mm of perp] (lob) {Libro de órdenes\\del contrato\\(\emph{bid}, \emph{ask}, \emph{spread})};
\node[tgt, right=22mm of lob] (mid) {Precio del contrato\\$m_t \in [0,1]$};
\node[box, below=14mm of mid, fill=gray!8] (pred) {Pregunta: ¿hacia dónde\\se mueve $m_t$ en los\\próximos segundos?};
\draw[->] (spot) -- (lob);
\draw[->] (perp) -- (lob);
\draw[->] (oracle) -- (lob);
\draw[->] (lob) -- (mid);
\draw[->] (mid) -- (pred);
\end{tikzpicture}
\caption{El entorno del problema. El precio del contrato binario se forma en su libro de
órdenes y está acoplado a referencias externas del mercado de Bitcoin. El objetivo es
anticipar su movimiento de muy corto plazo.}
\label{fig:mercado}
\end{figure}

\subsection{Definición del problema}

El precio del contrato no es el precio de Bitcoin: es una función no lineal y dependiente
del tiempo de la probabilidad de que el suceso se resuelva afirmativamente. Sin embargo,
está acoplado a las referencias externas observables que se acaban de citar. Un modelo
útil debe, por tanto, aprender dos cosas a la vez: la dinámica interna del libro de
órdenes y la semántica del contrato respecto a Bitcoin.

El problema que se aborda es de microestructura y de muy corto plazo: dado el estado del
mercado en un instante $t$ (libro de órdenes y referencias externas), predecir el
movimiento del precio medio del contrato en una ventana de pocas decenas de segundos. A
diferencia de la mayor parte de la literatura sobre mercados de predicción, que estudia
cómo se calibra el precio frente a la resolución final del evento, aquí el interés está en
los movimientos locales del precio, a escala de segundos.

Conviene fijar desde el principio un matiz que vertebra todo el trabajo: una buena
capacidad de \emph{predecir la dirección} no implica una estrategia rentable. Una métrica
de clasificación puede ser muy alta y, aun así, el sistema no ser explotable si opera con
retraso, si aprende artefactos del propio conjunto de datos o si ignora el coste real de
cruzar una orden. Por eso el problema no se evalúa por el acierto, sino por el resultado
económico neto, una vez descontados costes y latencia realistas.

\section{Motivación}

La motivación de este trabajo es doble. Desde el punto de vista aplicado, los mercados de
predicción de criptoactivos son un entorno joven, con menos competencia algorítmica que
los mercados financieros tradicionales, y con datos públicos accesibles; resulta natural
preguntarse si contienen señales explotables. Desde el punto de vista metodológico —que es
el que da peso a una memoria de máster— el entorno es un excelente banco de pruebas para
una cuestión central en el aprendizaje automático aplicado a finanzas: la distancia entre
\emph{predecir bien} y \emph{ganar dinero}.

Esa distancia es el hilo conductor. Como se mostrará, un clasificador direccional sencillo
alcanza una precisión cercana al 90\,\% y, pese a ello, deja de ser rentable en cuanto se
introduce una latencia de entrada realista de dos segundos. El coste y la ejecución, y no
el acierto, determinan el resultado. Asumir lo contrario es uno de los errores más comunes
y costosos en este campo, y documentarlo con evidencia cuantitativa es, en sí mismo, una
aportación.

\section{Objetivos}

El objetivo general es \textbf{construir y validar, con un protocolo honesto, una
metodología completa para evaluar si el aprendizaje automático y profundo pueden anticipar
el movimiento de corto plazo del precio de los contratos de Bitcoin de Polymarket bajo
costes y latencia realistas}. De él se derivan los siguientes objetivos específicos:

\begin{enumerate}[noitemsep]
  \item Diseñar un sistema de captura de datos multifuente, con control de calidad por
  sesión, que produzca un corpus alineado temporalmente y libre de fugas.
  \item Definir un protocolo de evaluación que combine validación estrictamente temporal,
  un modelo explícito de costes y una latencia medida empíricamente.
  \item Cuantificar el impacto de la latencia sobre una señal direccional fuerte, para
  delimitar dónde se rompe la cadena entre predicción y ejecución.
  \item Evaluar una progresión de modelos, de menor a mayor complejidad (líneas base
  tabulares, modelos sensibles a la latencia y arquitecturas profundas sobre el libro de
  órdenes), exigiendo que cada incremento de complejidad se justifique con evidencia.
  \item Determinar si existe un candidato que sobreviva a una validación fuera de muestra
  bajo coste, y delimitar con rigor sus límites y su incertidumbre.
\end{enumerate}

El presente trabajo se organiza como sigue. El Capítulo~\ref{ch::estado} revisa el estado
de la cuestión. El Capítulo~\ref{ch::metodos} describe los datos, las líneas base y la
metodología de modelado y entrenamiento. El Capítulo~\ref{ch::resultados} presenta y
analiza los resultados sobre el conjunto fuera de muestra. El Capítulo~\ref{ch::despliegue}
discute la estrategia de despliegue. El Capítulo~\ref{ch::etica} aborda los aspectos éticos
y las limitaciones. El Capítulo~\ref{ch::reproducibilidad} detalla la reproducibilidad, y
el Capítulo~\ref{ch::conclusiones} recoge las conclusiones. Por último, los anexos incluyen
un glosario de términos para el lector no familiarizado con las criptomonedas, el detalle de
las características y los hiperparámetros de los modelos.
